\section{Reading the same geometry in reverse}
\label{sec:inverse-question}

Part I asked a forward question: given electromagnetic parameters, which
current direction and magnitude are best?  The inverse question keeps every
operating constraint but changes what is fixed.  Let
\begin{equation}
 \gls{sym:theta}=(R_s,R_{\mathrm{exc}},L_d,L_q,L_{\mathrm{exc}},
 p,\psi_{\mathrm{pm}},\ldots)
 \label{eq:theta-definition}
\end{equation}
denote electromagnetic parameters, and reserve \(\mathcal I\) for the current
space containing \((i_d,i_q,i_{\mathrm{exc}})\).  The two spaces answer
different questions:
\begin{align}
 \boldsymbol\theta\ \hbox{fixed}:&\quad
 \text{which }\vect i\in\mathcal I\text{ is optimal?}\\
 \text{specification fixed}:&\quad
 \text{which }\boldsymbol\theta\text{ admits at least one feasible }\vect i?
\end{align}
Confusing these quantifiers is the most common conceptual error in an inverse
reading.  A current is not a machine parameter, and an electromagnetic
parameter vector is not yet a physical rotor and stator geometry.

The operating solver of Part I already supplies the required value functions.
For one requested torque-speed point define
\begin{equation}
 \gls{sym:thetafeas}=
 \left\{\boldsymbol\theta:\exists\vect i,
 \ M(\vect i;\boldsymbol\theta)=M_{\mathrm{req}},\ 
 P_{\mathrm{Cu}}(\vect i;\boldsymbol\theta)\leq P_{\mathrm{Cu},\max},\ 
 U(\vect i;\boldsymbol\theta,\omega_e)\leq U_{\max}\right\}.
 \label{eq:theta-feas-exists}
\end{equation}
Equivalently,
\begin{equation}
 \boxed{\boldsymbol\theta\in\Theta_{\mathrm{feas}}
 \iff M_{\max}(P_{\mathrm{Cu},\max},U_{\max},\omega_e;
 \boldsymbol\theta)\geq M_{\mathrm{req}}.}
 \label{eq:inverse-value-test}
\end{equation}
For a positive request the minimum-loss form is
\begin{equation}
 \boldsymbol\theta\in\Theta_{\mathrm{feas}}
 \iff P_{\min}(M_{\mathrm{req}},U_{\max},\omega_e;
 \boldsymbol\theta)\leq P_{\mathrm{Cu},\max}.
 \label{eq:inverse-pmin-test}
\end{equation}
Thus minimum loss and maximum torque are not separate design criteria.  They
are reciprocal level-set descriptions of the same feasibility boundary.

In current space, the boundary was the first tangency or intersection of a
torque surface with the loss and voltage sets.  In parameter space, the
boundary
\begin{equation}
 M_{\max}(P_{\mathrm{Cu},\max},U_{\max},\omega_e;
 \boldsymbol\theta)=M_{\mathrm{req}}
 \label{eq:design-capability-boundary}
\end{equation}
collects all machines for which that current-space contact has just appeared.
This is a projection of a higher-dimensional feasible set: the current is
existentially eliminated, not guessed or fixed in advance.

\paragraph{What can be inverted explicitly?}
When voltage is inactive and the model is linear, whitening provides closed
EESM and nonsalient-PSM inequalities.  With voltage, saliency, saturation, or
several operating points, a unique algebraic inverse should not be expected.
The honest output is then an implicit level surface such as
\cref{eq:design-capability-boundary}, a two-dimensional slice, or a monotone
root in one selected parameter.  This distinction prevents the phrase
``inverse design'' from hiding a second numerical optimization.

